\pdfbookmark{Общая характеристика работы}{characteristic}             % Закладка pdf
\section*{Общая характеристика работы}

\newcommand{\actuality}{\pdfbookmark[1]{Актуальность}{actuality}\underline{\textbf{\actualityTXT}}}
\newcommand{\progress}{\pdfbookmark[1]{Разработанность темы}{progress}\underline{\textbf{\progressTXT}}}
\newcommand{\aim}{\pdfbookmark[1]{Цели}{aim}\underline{{\textbf\aimTXT}}}
\newcommand{\tasks}{\pdfbookmark[1]{Задачи}{tasks}\underline{\textbf{\tasksTXT}}}
\newcommand{\aimtasks}{\pdfbookmark[1]{Цели и задачи}{aimtasks}\aimtasksTXT}
\newcommand{\object}{\pdfbookmark[1]{Объект исследования}{object}\underline{\textbf{\objectTXT}}}
\newcommand{\subject}{\pdfbookmark[1]{Предмет исследования}{subject}\underline{\textbf{\subjectTXT}}}
\newcommand{\novelty}{\pdfbookmark[1]{Научная новизна}{novelty}\underline{\textbf{\noveltyTXT}}}
\newcommand{\tinfluence}{\pdfbookmark[1]{Теоретическая значимость}{tinfluence}\underline{\textbf{\tinfluenceTXT}}}
\newcommand{\pinfluence}{\pdfbookmark[1]{Практическая значимость}{pinfluence}\underline{\textbf{\pinfluenceTXT}}}
\newcommand{\methods}{\pdfbookmark[1]{Методология и методы исследования}{methods}\underline{\textbf{\methodsTXT}}}
\newcommand{\defpositions}{\pdfbookmark[1]{Положения, выносимые на защиту}{defpositions}\underline{\textbf{\defpositionsTXT}}}
\newcommand{\passport}{\pdfbookmark[1]{Соответствие паспорту специальности}{passport}\underline{\textbf{\passportTXT}}}
\newcommand{\reliability}{\pdfbookmark[1]{Достоверность}{reliability}\underline{\textbf{\reliabilityTXT}}}
\newcommand{\probation}{\pdfbookmark[1]{Апробация}{probation}\underline{\textbf{\probationTXT}}}
\newcommand{\contribution}{\pdfbookmark[1]{Личный вклад}{contribution}\underline{\textbf{\contributionTXT}}}
\newcommand{\publications}{\pdfbookmark[1]{Публикации}{publications}\underline{\textbf{\publicationsTXT}}}

{\actuality} Декодирование по сигналам мозга~--- одна из центральных задач нейровизуализации и интерфейсов мозг--компьютер: по зарегистрированной активности требуется восстановить вызвавшее ее внешнее воздействие или его категорию. Регистрируемые сигналы представляют собой мультимодальные многомерные временные ряды. Функциональная магнитно-резонансная томография (фМРТ) дает высокое пространственное разрешение при низкой частоте сканирования и гемодинамическом лаге отклика, а электроэнцефалография (ЭЭГ)~--- высокое временное разрешение при слабой пространственной локализации. Эти модальности дополняют друг друга по пространству и времени, и их совместный анализ открывает возможности, недоступные каждой из них в отдельности.

Основное препятствие на этом пути~--- ограниченный объем выборки. Запись активности одного пациента сводится к одному протяженному временному ряду, число размеченных объектов измеряется десятками, тогда как размерность сигнала составляет десятки тысяч вокселей. В таких условиях нейросетевые модели, требующие больших обучающих выборок, теряют качество, а методы, не учитывающие структуру данных, не отделяют сигнал от шума. Поэтому актуальна разработка методов декодирования, которые явно используют временную структуру сигнала, прежде всего гемодинамический лаг, его пространственную структуру, то есть локализацию активных областей и ковариационные связи между ними, и межмодальную взаимодополняемость фМРТ и ЭЭГ, сохраняя при этом работоспособность при малом объеме выборки.

Отдельный вопрос составляет обоснованность применяемых конструкций. Проекция ковариационных матриц на касательное пространство риманова многообразия широко применяется в анализе нейросигналов, однако используется как эвристика: инвариантность признаков к преобразованию каналов, близость евклидовой геометрии касательного пространства к геометрии многообразия и устойчивость признаков к ошибке оценки ковариации по короткому ряду принимаются без явных констант. Получение таких гарантий необходимо для обоснованного выбора параметров метода и области его применимости.

{\aim} данной работы является разработка методов декодирования внешнего воздействия по мультимодальным многомерным временным рядам, использующих временную, пространственную и межмодальную структуру сигнала и сохраняющих работоспособность в условиях ограниченного объема выборки. Для достижения цели были поставлены и решены следующие {\tasks}:
\begin{enumerate}[beginpenalty=10000]
    \item Построить модель прогнозирования фМРТ-отклика по предъявленному стимулу и на ее основе метод оценки гемодинамического лага.
    \item Разработать метод классификации временных рядов фМРТ одного пациента на основе масок активности классов и проекции ковариационных матриц на касательное пространство риманова многообразия.
    \item Исследовать теоретические свойства касательных признаков ковариационных матриц: инвариантность к преобразованию каналов, искажение расстояний при проекции и устойчивость к шуму оценки ковариации.
    \item Построить мультимодальную архитектуру реконструкции стимула по одновременным записям фМРТ и ЭЭГ.
    \item Провести вычислительные эксперименты на открытых наборах данных и сравнить предложенные методы с нейросетевыми и классическими подходами.
\end{enumerate}

{\object} настоящей работы являются мультимодальные многомерные временные ряды, в частности ряды нейровизуализации, регистрируемые методами фМРТ и ЭЭГ в ответ на внешнее воздействие.

{\subject} данной работы являются методы декодирования внешнего воздействия, учитывающие временную, пространственную и межмодальную структуру таких рядов, и их теоретические свойства.

{\novelty}
\begin{enumerate}[beginpenalty=10000]
    \item Предложен метод оценки гемодинамического лага по минимуму ошибки авторегрессионного прогноза фМРТ-отклика по стимулу; показано, что в области наибольшего отклика ошибка имеет выраженный минимум, согласованный с физиологическими оценками лага.
    \item Предложен метод классификации временных рядов фМРТ, в котором объект описывается отдельно с точки зрения каждого класса: маска активности класса отбирает воксели по кросс-корреляции со стимулом, а ковариация их рядов проецируется на касательное пространство в точке риманова среднего; размерность признаков не зависит от объема выборки.
    \item Доказаны инвариантность касательных признаков к невырожденному преобразованию каналов, двусторонняя оценка искажения расстояний с неулучшаемой константой, выраженной через спектральный разброс касательных векторов, и оценка концентрации ошибки оценки ковариации в геодезической метрике, распределение которой не зависит от истинной ковариации.
    \item Предложена мультимодальная архитектура реконструкции стимула по одновременным записям фМРТ и ЭЭГ с контрастивным выравниванием совместных представлений с пространством CLIP и двухстадийной диффузионной генерацией; показано, что объединение модальностей повышает качество реконструкции по сравнению с каждой из них в отдельности.
\end{enumerate}

{\tinfluence} полученных результатов заключается в строгом обосновании широко применяемой конструкции касательных признаков ковариационных матриц: доказаны инвариантность признаков к преобразованию каналов, двусторонняя оценка искажения расстояний с явной константой и оценка шума оценки ковариации, связывающая число каналов и длину ряда. Результаты не зависят от природы сигнала и применимы к любым многоканальным временным рядам, описываемым ковариацией каналов.

{\pinfluence} настоящей работы заключается в применимости методов оценки лага и классификации к выборке одного пациента объемом в десятки объектов, где нейросетевые модели не обучаются, и в межсубъектной нейросетевой архитектуре реконструкции стимула, обучаемой на одновременных записях фМРТ и ЭЭГ многих пациентов; методы проверены на открытых наборах данных по воспроизводимым протоколам, а их программные реализации опубликованы в открытых репозиториях. Оценка лага используется при построении масок активности и при синхронизации модальностей, а метод восстановления отсутствующих каналов ЭЭГ позволяет применять обученные модели к записям с неполным монтажом.

{\methods} Для достижения поставленных в диссертационной работе целей используются:
\begin{enumerate}[beginpenalty=10000]
    \item Методы линейной алгебры и римановой геометрии: аффинно-инвариантная метрика на многообразии симметричных положительно определенных матриц, матричные функции, дифференциал матричной экспоненты.
    \item Методы теории вероятностей и математической статистики: неасимптотические оценки сингулярных чисел гауссовых матриц, проверка статистических гипотез с поправкой на множественные сравнения, доверительные интервалы.
    \item Методы машинного обучения и вычислительного эксперимента: регуляризованная линейная регрессия, логистическая регрессия, контрастивное обучение представлений, диффузионные генеративные модели, воспроизводимая постановка экспериментов на открытых данных.
\end{enumerate}

{\defpositions}
\begin{enumerate}[beginpenalty=10000]
    \item Линейная авторегрессионная модель прогнозирования фМРТ-отклика по предъявленному стимулу и метод оценки гемодинамического лага BOLD-сигнала на ее основе.
    \item Метод классификации временных рядов фМРТ на основе масок активности классов и проекции ковариационных матриц на касательное пространство риманова многообразия, устойчивый в условиях ограниченного объема выборки одного пациента.
    \item Теоретические свойства проекции ковариационных матриц на касательное пространство в точке риманова среднего: инвариантность признаков к невырожденному преобразованию каналов, двусторонняя оценка искажения расстояний с неулучшаемой константой и оценка ошибки оценки ковариации по длине ряда.
    \item Мультимодальная архитектура реконструкции стимула по одновременным сигналам фМРТ и ЭЭГ с контрастивным выравниванием совместных представлений и двухстадийной диффузионной генерацией.
\end{enumerate}

{\passport} Диссертационная работа соответствует паспорту научной специальности 1.2.1~--- <<Искусственный интеллект и машинное обучение>>, а именно в наибольшей степени следующим направлениям исследований: п.~12 <<Исследования в области надежности, устойчивости и переобучения систем класса ИИ>>; п.~15 <<Математические исследования в области статистики, логики, алгебры, топологии, анализа функций и других областях, ориентированные на решение задач искусственного интеллекта и машинного обучения>>; п.~16 <<Исследования в области специальных методов оптимизации, проблем сложности и снижения размерности>>; п.~17 <<Исследования в области многослойных алгоритмических конструкций, в том числе многослойных нейросетей>>.

{\reliability} полученных результатов обеспечивается корректностью применения апробированного математического аппарата римановой геометрии, теории вероятностей и математической статистики, математически строгими доказательствами теоретических утверждений, а также экспериментальной проверкой разработанных методов на открытых наборах данных с фиксированными протоколами разбиения, статистическими критериями с поправкой на множественные сравнения и доверительными интервалами. Для вычислительных экспериментов даются детальные описания и опубликован программный код, обеспечивающие их воспроизводимость. Результаты опубликованы в рецензируемых научных изданиях и доложены на российских и международных конференциях по машинному обучению и анализу данных.

{\probation} Основные результаты работы докладывались и обсуждались на следующих научных конференциях:
\begin{enumerate}
    %% \item Международная конференция <<Интеллектуализация обработки информации (ИОИ-2026)>>, сентябрь 2026~г.
    \item 68-я Всероссийская научная конференция МФТИ,
    Долгопрудный, 30 марта -- 4 апреля 2026~г.
    \item X Международная конференция по искусственному интеллекту AI Journey 2025,
    Москва, 19--21 ноября 2025~г.
    \item 22-я Всероссийская конференция с международным участием
    <<Математические методы распознавания образов (ММРО-2025)>>,
    Муром, 22--26 сентября 2025~г.
    \item 67-я Всероссийская научная конференция МФТИ,
    Долгопрудный, 31 марта -- 5 апреля 2025~г.
    \item XXVI Международная научно-техническая конференция <<Нейроинформатика-2024>>,
    Долгопрудный, 21--25 октября 2024~г.
    \item 66-я Всероссийская научная конференция МФТИ,
    Долгопрудный, 1--6 апреля 2024~г.
\end{enumerate}


{\contribution} Все приведенные результаты, кроме отдельно оговоренных случаев, получены диссертантом лично при научном руководстве к.ф.-м.н.~А.\,В.~Грабового.

\ifnumequal{\value{bibliosel}}{0}
{%%% Встроенная реализация с загрузкой файла через движок bibtex8. (При желании, внутри можно использовать обычные ссылки, наподобие `\cite{vakbib1,vakbib2}`).
    {\publications} Основные результаты по теме диссертации изложены
    в~XX~печатных изданиях,
    X из которых изданы в журналах, рекомендованных ВАК,
    X "--- в тезисах докладов.
}%
{%%% Реализация пакетом biblatex через движок biber
    \begin{refsection}[bl-author, bl-registered]
        % Это refsection=1.
        % Процитированные здесь работы:
        %  * подсчитываются, для автоматического составления фразы "Основные результаты ..."
        %  * попадают в авторскую библиографию, при usefootcite==0 и стиле `\insertbiblioauthor` или `\insertbiblioauthorgrouped`
        %  * нумеруются там в зависимости от порядка команд `\printbibliography` в этом разделе.
        %  * при использовании `\insertbiblioauthorgrouped`, порядок команд `\printbibliography` в нем должен быть тем же (см. biblio/biblatex.tex)
        %
        % Невидимый библиографический список для подсчета количества публикаций:
        \phantom{\printbibliography[heading=nobibheading, section=1, env=countauthorvak,          keyword=biblioauthorvak]%
        \printbibliography[heading=nobibheading, section=1, env=countauthorwos,          keyword=biblioauthorwos]%
        \printbibliography[heading=nobibheading, section=1, env=countauthorscopus,       keyword=biblioauthorscopus]%
        \printbibliography[heading=nobibheading, section=1, env=countauthorconf,         keyword=biblioauthorconf]%
        \printbibliography[heading=nobibheading, section=1, env=countauthorother,        keyword=biblioauthorother]%
        \printbibliography[heading=nobibheading, section=1, env=countauthor,             keyword=biblioauthor]%
        \printbibliography[heading=nobibheading, section=1, env=countauthorvakscopuswos, filter=vakscopuswos]%
        \printbibliography[heading=nobibheading, section=1, env=countauthorscopuswos,    filter=scopuswos]}%
        %
        \nocite{*}%
        %
        {\publications} \ifsynopsis\textit{Список публикаций приведен в конце автореферата. }\else\textit{Список публикаций приведен в конце диссертации. }\fi Основные результаты по теме диссертации изложены в~\arabic{citeauthor}~научных работах, из которых~\arabic{citeauthorscopuswos} "--- в~периодических научных журналах, индексируемых Web of~Science и Scopus, а остальные~\arabic{citeauthorconf} "--- в~тезисах докладов российских и международных конференций.
    \end{refsection}%
    \begin{refsection}[bl-author, bl-registered]
        % Это refsection=2.
        % Процитированные здесь работы:
        %  * попадают в авторскую библиографию, при usefootcite==0 и стиле `\insertbiblioauthorimportant`.
        %  * ни на что не влияют в противном случае
        \nocite{dorin2024forecasting}
        \nocite{dorin2025enhancing}
        \nocite{dorin2026decoding}
        \nocite{dorin2025mmro}
        \nocite{dorin2024conf66mipt}
        \nocite{dorin2025conf67mipt}
        \nocite{dorin2026conf68mipt}
        \nocite{dorin2025aij}
    \end{refsection}%
}
 % Характеристика работы по структуре во введении и в автореферате не отличается (ГОСТ Р 7.0.11, пункты 5.3.1 и 9.2.1), потому её загружаем из одного и того же внешнего файла, предварительно задав форму выделения некоторым параметрам

\pdfbookmark[1]{Объем и структура работы}{structure}\underline{\textbf{Объем и структура работы.}} Диссертация состоит из введения, четырех глав, заключения, списка сокращений и условных обозначений и списка литературы. Полный объем диссертации составляет 121~страницу, включая 31~рисунок и 14~таблиц. Список литературы содержит 129~наименований.

\pdfbookmark{Основное содержание работы}{description}                 % Закладка pdf
\section*{Основное содержание работы}

\textbf{Во введении} обоснована актуальность темы, сформулированы цель и задачи исследования, научная новизна, теоретическая и практическая значимость полученных результатов, положения, выносимые на защиту, приведены сведения о достоверности и апробации результатов и о публикациях автора.

\textbf{В первой главе} выполнен анализ предметной области и дана общая постановка задачи. Рассмотрены применения интерфейсов мозг--компьютер: управление протезами руки и экзоскелетами, восстановление речи и зрения, диагностика нарушений сознания и болезни Альцгеймера. Рассмотрены также свойства модальностей фМРТ и ЭЭГ: первая дает высокое пространственное разрешение при низкой частоте сканирования и гемодинамическом лаге сигнала, зависящего от уровня оксигенации крови (англ. Blood-Oxygen-Level-Dependent, BOLD), вторая~--- высокое временное разрешение при слабой пространственной локализации. Описаны парадигмы стимуляции и способы выделения активных областей мозга, применение римановой геометрии ковариационных матриц к анализу нейросигналов, одномодальные методы декодирования и реконструкции визуальной информации на основе контрастивного обучения и диффузионных генеративных моделей, визуальные энкодеры и открытые наборы данных, в том числе наборы с одновременной регистрацией фМРТ и ЭЭГ. Показано, что существующие нейросетевые методы требуют больших обучающих выборок, методы реконструкции опираются на одну модальность, а свойства касательных признаков ковариационных матриц используются без явных гарантий.

Декодирование рассматривается как единая задача. Пусть $\mathcal{X}_{(m)}$~--- пространство наблюдений модальности $m = 1, \ldots, M$, $\mathcal{X} = \mathcal{X}_{(1)} \times \cdots \times \mathcal{X}_{(M)}$~--- совместное пространство данных, а $\mathcal{Y}$~--- пространство стимулов или их характеристик. Требуется построить отображение $\mathbf{g}\colon \mathcal{X} \to \mathcal{Y}$. В работе последовательно решаются три задачи, упорядоченные по сложности описываемой структуры данных и показанные на Рисунке~\ref{fig:syn_overview}: оценка гемодинамического лага по временной динамике фМРТ, классификация временных рядов фМРТ с учетом их пространственной структуры и мультимодальная реконструкция стимула по одновременным записям фМРТ и ЭЭГ при $M = 2$.

\begin{figure}[h!]
    \centering
    \includegraphics[width=\textwidth]{scheme_overview.pdf}
    \caption{Три постановки задачи декодирования. Штриховые линии~--- учет гемодинамического лага $\Delta t$, необходимость которого показана в первой постановке, при построении масок активности и при синхронизации модальностей}
    \label{fig:syn_overview}
\end{figure}

\textbf{Во второй главе} рассматривается первая постановка, учитывающая только временную структуру сигнала~\cite{dorin2024forecasting}. Сигнал фМРТ возникает с гемодинамическим лагом $\Delta t$ относительно момента предъявления стимула. Для оценки лага решается обратная к декодированию задача: по стимулу $\mathbf{y}(t)$ прогнозируется скан $\mathbf{F}(t) \in \mathbb{R}^{X \times Y \times Z}$, $\mathbf{F}(t) \approx \mathbf{q}[\mathbf{y}(t - \Delta t)]$, а искомый лаг есть значение $\Delta t$, при котором ошибка прогноза минимальна. Отображение $\mathbf{q} = \boldsymbol{\varphi} \circ \boldsymbol{\psi}$ составлено из фиксированной векторизации кадра предобученным энкодером, $\mathbf{x}(t) = \boldsymbol{\psi}[\mathbf{y}(t - \Delta t)] \in \mathbb{R}^{d}$, и обучаемого линейного отображения $\boldsymbol{\varphi}$. Последовательность сканов считается марковской, поэтому прогнозируется приращение $\boldsymbol{\delta}(t) = \mathbf{F}(t) - \mathbf{F}(t-1)$, для каждого вокселя~--- линейной регрессией с $L_2$-регуляризацией:
\begin{equation}\label{eq:syn_lag}
\widehat{\mathbf{w}}_{xyz} = \arg\min_{\mathbf{w}} \sum_{t=2}^{\mathcal{T}} \big( \langle \mathbf{x}(t), \mathbf{w} \rangle - \delta_{xyz}(t) \big)^2 + \lambda \|\mathbf{w}\|_2^2
= (\mathbf{X}^{\top}\mathbf{X} + \lambda\mathbb{I})^{-1} \mathbf{X}^{\top} \boldsymbol{\Delta}_{xyz},
\end{equation}
где $\mathbf{X}$~--- матрица эмбеддингов кадров, $\boldsymbol{\Delta}_{xyz}$~--- вектор приращений вокселя. Веса вокселей объединяются в матрицу $\widehat{\mathbf{W}}$, и последовательность сканов восстанавливается авторегрессионно: $\operatorname{vec}\widehat{\mathbf{F}}(t) = \operatorname{vec}\mathbf{F}(t-1) + \widehat{\mathbf{W}}\mathbf{x}(t)$.

Эксперимент проведен на открытом наборе данных с фМРТ-записями 30 пациентов, просматривавших один и тот же короткий фильм: длительность записи 390~с, частота кадров 25~Гц, частота сканирования 1,64~Гц. Кадры кодируются сетью ResNet152, сканы сжимаются трехмерным пулингом и нормируются. По всему объему скана зависимость среднеквадратичной ошибки (англ. Mean Squared Error, MSE) от $\Delta t$ не имеет выраженного минимума из-за шума областей, не связанных со зрением; на локализованной затылочной доле минимум отчетлив и достигается при $\Delta t \approx 5$~с, что согласуется с физиологическими оценками лага, как показано на Рисунке~\ref{fig:syn_mse_dt}. Показано также, что оптимальный коэффициент регуляризации $\lambda \approx 1000$ не зависит от степени сжатия сканов, что матрица весов одного пациента восстанавливает сканы другого практически с той же ошибкой и что замена эмбеддингов кадров неинформативными увеличивает ошибку в несколько раз. Эксперимент обосновывает явный учет лага в последующих методах. Он показывает также, что без выделения информативной области зависимость качества от лага не проявляется.

\begin{figure}[h!]
    \centering
    \includegraphics[width=\textwidth]{figs/mse_dt.pdf}
    \caption{Зависимость MSE прогноза от гемодинамического лага $\Delta t$: по всему скану, слева, и на локализованной затылочной доле, справа, где минимум достигается при $\Delta t \approx 5$~с}
    \label{fig:syn_mse_dt}
\end{figure}

\textbf{В третьей главе} рассматривается вторая постановка, дополняющая временное описание пространственным: временные ряды фМРТ одного пациента классифицируются по категориям стимулов~\cite{dorin2025enhancing}. Объект~--- ряд $\mathbf{S} = [\mathbf{F}(1), \ldots, \mathbf{F}(\mathcal{T})]$ из $\mathcal{T}$ сканов с меткой $y \in \{1, \ldots, L\}$; выборка $\mathfrak{D} = \{(\mathbf{S}_i, y_i)\}_{i=1}^{N}$ одного пациента содержит десятки объектов. Предложенная модель есть суперпозиция $\mathbf{h} = \boldsymbol{\eta} \circ \boldsymbol{\psi} \circ \mathcal{A}$, где $\mathcal{A}$~--- пространственное сжатие сканов, $\boldsymbol{\psi} = \boldsymbol{\psi}_1 \oplus \cdots \oplus \boldsymbol{\psi}_L$~--- кодировщик, описывающий объект отдельно с точки зрения каждого класса, $\boldsymbol{\eta}$~--- классификатор. Схема модели приведена на Рисунке~\ref{fig:syn_cls}.

\begin{figure}[h!]
    \centering
    \includegraphics[width=\textwidth]{cls/cls_method_scheme.pdf}
    \caption{Схема модели классификации. Маска активности $\mathcal{M}^{y}$ класса $y$ отбирает ряды активных вокселей $\mathbf{Y}_y$, их ковариационная матрица $\mathbf{R}_y$ проецируется на касательное пространство в точке риманова среднего $\bar{\mathbf{R}}_y$, что дает вектор $\mathbf{p}_y$ соответствия объекта классу; векторы всех классов конкатенируются в признак $\mathbf{p}$}
    \label{fig:syn_cls}
\end{figure}

Маска активности класса $y$ строится по кросс-корреляции сигнала вокселя с бинарным рядом стимулов $\mathbf{u}$ этой категории при гемодинамическом лаге:
\begin{equation*}
\rho_{xyz} = \frac{1}{\mathcal{T}-1} \sum_{t} \mathring{u}(t)\, \mathring{F}_{xyz}(t + \Delta t),
\end{equation*}
где кружок обозначает Z-нормализацию ряда. Ряды и стимулы обучающих объектов класса конкатенируются по времени, и в маску $\mathcal{M}^{y}$ входят $h_y$ вокселей с наибольшей корреляцией; значимость корреляции проверяется критерием Кендалла с поправкой Холма на множественную проверку. Отображение $\boldsymbol{\psi}_y$ применяет маску и оставляет матрицу рядов активных вокселей $\mathbf{Y} \in \mathbb{R}^{h \times \mathcal{T}}$, по которой вычисляется ковариационная матрица $\mathbf{R} = (\mathcal{T}-1)^{-1} \mathbf{Y}\mathbf{Y}^{\top}$~--- точка многообразия $\mathbb{S}_{+}^{h}$ симметричных положительно определенных матриц с аффинно-инвариантной метрикой и геодезическим расстоянием $\delta_{\mathcal{R}}(\mathbf{R}_1, \mathbf{R}_2) = \|\log(\mathbf{R}_1^{-1/2}\mathbf{R}_2\mathbf{R}_1^{-1/2})\|_F$. Ковариации объектов проецируются на касательное пространство в точке риманова среднего обучающей выборки $\bar{\mathbf{R}} = \arg\min_{\mathbf{R}} \sum_i \delta_{\mathcal{R}}^2(\mathbf{R}, \mathbf{R}_i)$, что составляет отображение TSM (англ. Tangent Space Mapping):
\begin{equation}\label{eq:syn_feat}
\mathbf{S}_i = \log\big( \bar{\mathbf{R}}^{-1/2} \mathbf{R}_i \bar{\mathbf{R}}^{-1/2} \big), \qquad
\mathbf{p}_i = \operatorname{vec}_{\sqrt{2}}(\mathbf{S}_i) \in \mathbb{R}^{d}, \qquad d = \frac{h(h+1)}{2},
\end{equation}
где $\operatorname{vec}_{\sqrt{2}}$ записывает верхний треугольник симметричной матрицы в вектор с множителем $\sqrt{2}$ у внедиагональных элементов, так что $\|\mathbf{p}_i\|_2 = \|\mathbf{S}_i\|_F$. Признаки всех классов конкатенируются; их размерность определяется числами $h_y$ и не зависит от объема выборки. Для формирования выборки из одной протяженной записи предложен метод сегментации, выделяющий сегменты фиксированной длины с преобладающей категорией стимула.

Для конструкции~\eqref{eq:syn_feat} доказаны следующие утверждения. Пусть $\omega(\mathbf{S}) = \lambda_{\max}(\mathbf{S}) - \lambda_{\min}(\mathbf{S})$~--- спектральный разброс симметричной матрицы и $g(x) = \sinh(x/2)/(x/2)$.

\begin{proposition}\label{prop:syn_affine}
Пусть $\mathbf{A} \in \mathbb{R}^{h \times h}$ обратима и $\mathbf{R}_i' = \mathbf{A}\mathbf{R}_i\mathbf{A}^{\top}$, $i = 1, \ldots, N$. Тогда $\bar{\mathbf{R}}' = \mathbf{A}\bar{\mathbf{R}}\mathbf{A}^{\top}$; существует ортогональная матрица $\mathbf{O} \in \mathbb{R}^{d \times d}$, одна и та же для всех объектов, такая что $\mathbf{p}_i' = \mathbf{O}\mathbf{p}_i$; матрица весов решения $L_2$-регуляризованной логистической регрессии на признаках $\{\mathbf{p}_i'\}$ равна $\hat{\mathbf{W}}\mathbf{O}^{\top}$, где $\hat{\mathbf{W}}$~--- матрица весов решения на $\{\mathbf{p}_i\}$, и предсказания на любом объекте совпадают.
\end{proposition}

Следовательно, результат классификации не зависит от единиц измерения сигнала, масштабирования и невырожденного смешивания отобранных вокселей.

\begin{theorem}\label{thm:syn_isometry}
Для признаков~\eqref{eq:syn_feat} справедливы следующие утверждения: (а)~$\|\mathbf{p}_i\|_2 = \delta_{\mathcal{R}}(\bar{\mathbf{R}}, \mathbf{R}_i)$; (б)~$\|\mathbf{p}_i - \mathbf{p}_j\|_2 \leqslant \delta_{\mathcal{R}}(\mathbf{R}_i, \mathbf{R}_j)$ для любых $i, j$; (в)~если $\omega(\mathbf{S}_i) \leqslant \omega$ для всех $i$, то $\delta_{\mathcal{R}}(\mathbf{R}_i, \mathbf{R}_j) \leqslant g(\omega)\,\|\mathbf{p}_i - \mathbf{p}_j\|_2$; (г)~если $\delta_{\mathcal{R}}(\bar{\mathbf{R}}, \mathbf{R}_i) \leqslant \rho$ для всех $i$, то $\delta_{\mathcal{R}}(\mathbf{R}_i, \mathbf{R}_j) \leqslant c(\rho)\,\|\mathbf{p}_i - \mathbf{p}_j\|_2$, где $c(\rho) = g(\sqrt{2}\rho)$.
\end{theorem}

\begin{proposition}\label{prop:syn_sharp}
Пусть $\omega > 0$. Для любой пары различных объектов неравенства пунктов~(в) и~(г) теоремы~\ref{thm:syn_isometry} строгие. Пусть $\mathbf{S}(\theta) = \mathbf{Q}_{\theta} \operatorname{diag}(\omega/2, -\omega/2, 0, \ldots, 0) \mathbf{Q}_{\theta}^{\top}$, где $\mathbf{Q}_{\theta}$~--- поворот на угол $\theta$ в плоскости первых двух базисных векторов. Тогда $\omega(\mathbf{S}(\theta)) = \omega$ для всех $\theta$ и
\begin{equation*}
\lim_{\theta \to 0} \frac{\delta_{\mathcal{R}}\big(\exp \mathbf{S}(-\theta/2), \exp \mathbf{S}(\theta/2)\big)}{\|\mathbf{S}(\theta/2) - \mathbf{S}(-\theta/2)\|_F} = g(\omega),
\end{equation*}
то есть константы теоремы~\ref{thm:syn_isometry} неулучшаемы, но достигаются лишь в пределе.
\end{proposition}

Искажение расстояний порождается не удаленностью объектов от опорной точки, а некоммутативностью их касательных векторов, и управляется спектральным разбросом $\omega$, равным логарифму числа обусловленности отбеленной ковариации $\bar{\mathbf{R}}^{-1/2}\mathbf{R}_i\bar{\mathbf{R}}^{-1/2}$.

\begin{theorem}\label{thm:syn_conc}
Пусть столбцы матрицы $\mathbf{Y} \in \mathbb{R}^{h \times \mathcal{T}}$~--- независимые случайные векторы с распределением $\mathcal{N}(\mathbf{0}, \mathbf{R})$, $\mathbf{R} \in \mathbb{S}_{+}^{h}$, и $\hat{\mathbf{R}} = \mathcal{T}^{-1}\mathbf{Y}\mathbf{Y}^{\top}$. Тогда (а)~распределение случайной величины $\delta_{\mathcal{R}}(\hat{\mathbf{R}}, \mathbf{R})$ не зависит от $\mathbf{R}$, а только от $h$ и $\mathcal{T}$; (б)~для любого $t > 0$ такого, что $\varepsilon = \sqrt{h/\mathcal{T}} + t < 1$, с вероятностью не меньше $1 - 2\exp(-\mathcal{T}t^2/2)$
\begin{equation*}
\delta_{\mathcal{R}}(\hat{\mathbf{R}}, \mathbf{R}) \leqslant \frac{2\sqrt{h}\,\varepsilon}{1 - \varepsilon}.
\end{equation*}
\end{theorem}

По пункту~(б) теоремы~\ref{thm:syn_isometry} ошибка признака, вычисленного по оценке ковариации, не превосходит $\delta_{\mathcal{R}}(\hat{\mathbf{R}}_i, \mathbf{R}_i)$, поэтому шум оценки растет как $h/\sqrt{\mathcal{T}}$. Это обосновывает отбор в маску небольшого числа $h \ll \mathcal{T}$ наиболее коррелированных со стимулом вокселей.

Вычислительный эксперимент выполнен на открытом наборе данных с записями фМРТ при пассивном просмотре изображений восьми категорий, снабженном масками активности, размеченными нейроучеными. Из записи каждого пациента сегментацией получено $N = 96$ рядов длины $\mathcal{T} = 19$, по 12 на категорию. Маски алгоритма состоят из статистически значимых вокселей и покрывают разметку нейроученых, а на случайном ряде стимулов значимых вокселей не содержат. Результаты классификации при $h_y = 10$ приведены в Таблице~\ref{tab:syn_cls}: качество монотонно растет с долей обучающей выборки, а исключение любой из компонент модели, как и переход к нейросетевым кодировщикам на основе сети долгой краткосрочной памяти (англ. Long Short-Term Memory, LSTM) или механизма внимания, снижает точность в несколько раз. Все различия с предложенной моделью статистически значимы по t-критерию для связанных выборок с поправкой Холма.

\begin{table}[h!]
    \centering
    \small
    \caption{Качество классификации на тестовой выборке, среднее и стандартное отклонение по пациентам}\label{tab:syn_cls}
    \begin{tabular}{lcc}
        \toprule
        Модель & Accuracy & Macro F1-score \\
        \midrule
        Предложенная модель, логистическая регрессия & $0{,}61 \pm 0{,}04$ & $0{,}56 \pm 0{,}03$ \\
        Предложенная модель, перцептрон & $0{,}64 \pm 0{,}05$ & $0{,}60 \pm 0{,}05$ \\
        Без проекции на касательное пространство & $0{,}24 \pm 0{,}11$ & $0{,}22 \pm 0{,}11$ \\
        Без масок классов & $0{,}25 \pm 0{,}13$ & $0{,}23 \pm 0{,}13$ \\
        Маски нейроученых и TSM & $0{,}10 \pm 0{,}05$ & $0{,}09 \pm 0{,}04$ \\
        Маски нейроученых и LSTM & $0{,}12 \pm 0{,}02$ & $0{,}03 \pm 0{,}01$ \\
        Маски нейроученых и механизм внимания & $0{,}14 \pm 0{,}04$ & $0{,}04 \pm 0{,}03$ \\
        \bottomrule
    \end{tabular}
\end{table}

Теоретические оценки проверены на синтетических данных с явно заданными ковариациями. Предельный переход предложения~\ref{prop:syn_sharp} воспроизводится численно. С ростом геодезического радиуса облака объектов точность TSM совпадает с точностью классификатора по минимальному расстоянию до среднего, работающего на многообразии, и превосходит проекцию в единичной матрице и евклидову векторизацию ковариаций, тогда как фактическое искажение расстояний остается много меньше границ теоремы~\ref{thm:syn_isometry}. При замене истинных ковариаций выборочными ошибка убывает как $\mathcal{T}^{-1/2}$ и пропорциональна $h$, а граница теоремы~\ref{thm:syn_conc} убывает с той же скоростью.

\textbf{В четвертой главе} рассматривается третья постановка, расширяющая пространственно-временное описание на вторую модальность: по одновременным сигналам фМРТ и ЭЭГ реконструируются визуальные стимулы~\cite{dorin2026decoding, dorin2025mmro}. Задан набор триплетов $\mathfrak{D} = \{(\mathbf{F}_i, \mathbf{E}_i, \mathbf{I}_i)\}_{i=1}^{N}$, где $\mathbf{F}_i$~--- фМРТ-скан, $\mathbf{E}_i \in \mathbb{R}^{K \times T}$~--- окно записи ЭЭГ с $K$ каналами, $\mathbf{I}_i$~--- изображение; требуется построить отображение $\mathbf{f}\colon \mathcal{X}_f \times \mathcal{X}_e \to \mathcal{Y}$ в пространство изображений. Метод состоит из двух этапов, показанных на Рисунке~\ref{fig:syn_bimodal}: контрастивного выравнивания совместных представлений фМРТ и ЭЭГ с эмбеддингами изображений модели CLIP (англ. Contrastive Language-Image Pre-training) и двухстадийной диффузионной генерации.

\begin{figure}[h!]
    \centering
    \includegraphics[width=\textwidth]{figs/figure_1.pdf}
    \caption{Схема реконструкции изображений по сигналам фМРТ и ЭЭГ: контрастивное обучение энкодера сигналов мозга с замороженными эмбеддингами CLIP, затем диффузионная априорная модель в пространстве CLIP и предобученная диффузионная модель генерации изображений}
    \label{fig:syn_bimodal}
\end{figure}

На этапе предобработки каждому скану с учетом лага $\Delta t = 5$~с сопоставляются окно ЭЭГ и кадр видеостимула, для фМРТ по накопленной во времени вариации сигнала и порогу Оцу строится индивидуальная маска информативных вокселей, а записи ЭЭГ приводятся к полному набору каналов, поскольку набор зарегистрированных электродов меняется от записи к записи. Отсутствующий канал восстанавливается взвешенным средним $n$ ближайших по геометрии монтажа электродов с весами, обратными геодезическому расстоянию по поверхности скальпа. Предложено вероятностное обобщение этой интерполяции на основе согласования потоков (англ. flow matching). Пусть $\mathbf{M} \in \{0,1\}^{K \times T}$~--- маска наблюдаемости, $\bar{\mathbf{M}} = \mathbf{1} - \mathbf{M}$, $\mathbf{E}_{\mathrm{obs}} = \mathbf{M} \odot \mathbf{E}$, $\mathbf{D}$~--- матрица расстояний между электродами, $\boldsymbol{\varepsilon}$~--- гауссовский шум. Наблюдаемые значения неподвижны, а пропущенные движутся от шума к данным:
\begin{equation*}
\mathbf{E}_t = \mathbf{M} \odot \mathbf{E} + \bar{\mathbf{M}} \odot \big( (1-t)\,\boldsymbol{\varepsilon} + t\,\mathbf{E} \big), \qquad t \in [0,1].
\end{equation*}
Векторное поле $v_{\theta}(\mathbf{E}_t, t, \mathbf{E}_{\mathrm{obs}}, \mathbf{M}; \mathbf{D})$ задается нейросетью, чередующей свертки по времени и внимание по каналам со смещением логитов, зависящим от расстояния между электродами, и обучается регрессией по пропущенным позициям:
\begin{equation}\label{eq:syn_imp}
\mathcal{L}_{\mathrm{imp}}(\theta) = \mathbb{E} \Big\| \bar{\mathbf{M}} \odot \big( v_{\theta}(\mathbf{E}_t, t, \mathbf{E}_{\mathrm{obs}}, \mathbf{M}; \mathbf{D}) - (\mathbf{E} - \boldsymbol{\varepsilon}) \big) \Big\|_F^2 .
\end{equation}
Восстановление есть интегрирование уравнения $d\mathbf{E}_t/dt = \bar{\mathbf{M}} \odot v_{\theta}$ от $\mathbf{E}_0 = \mathbf{E}_{\mathrm{obs}} + \bar{\mathbf{M}} \odot \boldsymbol{\varepsilon}$ до $t = 1$; повторение с независимыми реализациями шума дает ансамбль восстановлений и оценку их неопределенности.

\begin{theorem}\label{thm:syn_imp}
Пусть $\boldsymbol{\xi}_1$~--- вектор пропущенных значений с распределением $p(\cdot \mid \mathbf{E}_{\mathrm{obs}}, \mathbf{M})$, $\mathbb{E}\|\boldsymbol{\xi}_1\|_2 < \infty$, $\boldsymbol{\xi}_0 \sim \mathcal{N}(\mathbf{0}, \mathbb{I})$ не зависит от $\boldsymbol{\xi}_1$, $\boldsymbol{\xi}_t = (1-t)\boldsymbol{\xi}_0 + t\boldsymbol{\xi}_1$, а поле $v^{*}(\boldsymbol{\xi}, t) = \mathbb{E}[\boldsymbol{\xi}_1 - \boldsymbol{\xi}_0 \mid \boldsymbol{\xi}_t = \boldsymbol{\xi}]$ непрерывно и липшицево по $\boldsymbol{\xi}$ равномерно по $t$. Тогда $v^{*}$ есть минимизатор функции потерь $\mathcal{L}_{\mathrm{imp}}$ по всем измеримым функциям состояния, времени и условия, а решение уравнения $\dot{\boldsymbol{\xi}} = v^{*}(\boldsymbol{\xi}, t)$ с начальным условием $\boldsymbol{\xi}(0) = \boldsymbol{\xi}_0$ существует, единственно и $\boldsymbol{\xi}(1) \sim p(\cdot \mid \mathbf{E}_{\mathrm{obs}}, \mathbf{M})$.
\end{theorem}

Энкодер ЭЭГ объединяет запись с обучаемым эмбеддингом пациента и применяет свертки по времени, по каналам и глобальную свертку с последующим многослойным перцептроном (англ. Multilayer Perceptron, MLP); энкодер фМРТ проецирует отобранные маской воксели индивидуальным для пациента линейным слоем в общее латентное пространство и обрабатывает их MLP с остаточными соединениями; модуль объединения конкатенирует эмбеддинги модальностей и преобразует их MLP в совместный эмбеддинг $\mathbf{Z}_c$ размерности пространства CLIP. Энкодеры обучаются на триплетах симметричной контрастивной функцией потерь по матрице логитов $\mathbf{L} = \mathbf{Z}_I^{n} (\mathbf{Z}_c^{n})^{\top} e^{\tau}$, где $\mathbf{Z}_I^{n}$ и $\mathbf{Z}_c^{n}$~--- нормированные эмбеддинги изображений замороженной модели CLIP и сигналов мозга, $\tau$~--- обучаемая температура:
\begin{equation}\label{eq:syn_contr}
\mathcal{L}_{\mathrm{contr}} = -\frac{1}{2B} \sum_{b=1}^{B} \left[ \log \frac{\exp(\mathbf{L}_{bb})}{\sum_{j} \exp(\mathbf{L}_{bj})} + \log \frac{\exp(\mathbf{L}_{bb})}{\sum_{i} \exp(\mathbf{L}_{ib})} \right].
\end{equation}
Реконструкция выполняется в две стадии. На первой в пространстве CLIP обучается диффузионная априорная модель $\boldsymbol{\epsilon}_{\boldsymbol{\theta}}(\mathbf{z}_I^{(t)}, t, \mathbf{z}_c)$ с функцией потерь $\mathbb{E}\|\boldsymbol{\epsilon} - \boldsymbol{\epsilon}_{\boldsymbol{\theta}}(\mathbf{z}_I^{(t)}, t, \mathbf{z}_c)\|_2^2$, уточняющая эмбеддинг сигналов мозга до согласованного с визуальным пространством эмбеддинга $\widehat{\mathbf{z}}_I$. На второй предобученная латентно-диффузионная модель SDXL (англ. Stable Diffusion XL), обусловленная на $\widehat{\mathbf{z}}_I$ через адаптер IP-Adapter (англ. Image Prompt Adapter), синтезирует изображение; энкодеры и генеративная модель при этом не дообучаются.

Эксперимент проведен на открытом наборе данных с одновременной регистрацией фМРТ и ЭЭГ у 22 пациентов при просмотре натуралистических видеороликов: 61 электрод ЭЭГ с частотой дискретизации 5000~Гц, частота сканирования фМРТ 0,476~Гц, около 15 тысяч триплетов, разделенных на обучающую и тестовую выборки в пропорции 80\% к 20\%. Качество измеряется метрикой CLIP-Score~--- косинусной близостью эмбеддинга реконструкции и совместного эмбеддинга сигналов мозга. Результаты в Таблице~\ref{tab:syn_clip} показывают, что объединение модальностей превосходит каждую из них в отдельности, а диффузионная априорная модель необходима для семантической согласованности реконструкций. На стимулах из домена, отсутствовавшего в обучении, качество реконструкции снижается, что ограничивает перенос метода без доменной адаптации.

\begin{table}[h!]
    \centering
    \small
    \caption{CLIP-Score реконструкций на тестовой выборке для трех видеостимулов, среднее по 22 пациентам; нижним индексом указаны стандартные отклонения}\label{tab:syn_clip}
    \resizebox{\textwidth}{!}{\begin{tabular}{lccc}
        \toprule
        Модель & Видео 1 & Видео 2 & Видео 3 \\
        \midrule
        ЭЭГ & $0{,}231_{\pm 0{,}004}$ & $0{,}262_{\pm 0{,}004}$ & $0{,}224_{\pm 0{,}005}$ \\
        фМРТ & $0{,}254_{\pm 0{,}004}$ & $0{,}270_{\pm 0{,}003}$ & $0{,}242_{\pm 0{,}003}$ \\
        фМРТ--ЭЭГ & $0{,}310_{\pm 0{,}003}$ & $0{,}312_{\pm 0{,}004}$ & $0{,}284_{\pm 0{,}003}$ \\
        ЭЭГ и априорная модель & $0{,}542_{\pm 0{,}005}$ & $0{,}557_{\pm 0{,}003}$ & $0{,}538_{\pm 0{,}003}$ \\
        фМРТ и априорная модель & $0{,}564_{\pm 0{,}005}$ & $0{,}571_{\pm 0{,}004}$ & $0{,}551_{\pm 0{,}005}$ \\
        фМРТ--ЭЭГ и априорная модель & $\mathbf{0{,}668}_{\pm 0{,}003}$ & $\mathbf{0{,}647}_{\pm 0{,}005}$ & $\mathbf{0{,}644}_{\pm 0{,}003}$ \\
        \bottomrule
    \end{tabular}}
\end{table}

\textbf{В заключении} сформулированы основные результаты работы, рекомендации по их использованию и перспективы дальнейшей разработки темы.

\pdfbookmark{Заключение}{conclusion}                                  % Закладка pdf
\section*{Заключение}
%% Согласно ГОСТ Р 7.0.11-2011:
%% 5.3.3 В заключении диссертации излагают итоги выполненного исследования, рекомендации, перспективы дальнейшей разработки темы.
%% 9.2.3 В заключении автореферата диссертации излагают итоги данного исследования, рекомендации и перспективы дальнейшей разработки темы.
В диссертационной работе решена задача разработки методов декодирования внешнего воздействия по мультимодальным многомерным временным рядам в условиях ограниченного объема выборки. Методы последовательно используют временную, пространственную и межмодальную структуру сигнала и проверены на данных нейровизуализации.

Основные результаты работы заключаются в следующем.
\begin{enumerate}
    \item Построена линейная авторегрессионная модель прогнозирования фМРТ-отклика по предъявленному стимулу и на ее основе предложен метод оценки гемодинамического лага по минимуму ошибки восстановления. В локализованной области отклика ошибка имеет выраженный минимум около пяти секунд, согласующийся с физиологическими оценками, тогда как по всему объему скана минимум не выражен. Этот результат обосновывает явный учет лага и выделение информативной области в последующих методах.
    \item Разработан метод классификации временных рядов фМРТ одного пациента: маски активности классов выделяются по кросс-корреляции сигнала со стимулом, а ковариация отобранных вокселей проецируется на касательное пространство в точке риманова среднего, так что объект описывается отдельно с точки зрения каждого класса. Размерность признаков не зависит от объема выборки; на данных с восемью категориями стимулов метод превосходит нейросетевые модели, а исключение любой из компонент значимо снижает качество.
    \item Доказаны теоретические свойства касательных признаков ковариационных матриц: инвариантность к невырожденному преобразованию каналов, не меняющая ответ логистической регрессии; двусторонняя оценка искажения расстояний с неулучшаемой константой, выраженной через спектральный разброс касательных векторов; оценка концентрации ошибки оценки ковариации в геодезической метрике, распределение которой не зависит от истинной ковариации.
    \item Предложена мультимодальная архитектура реконструкции стимула по одновременным записям фМРТ и ЭЭГ с контрастивным выравниванием совместных представлений с пространством CLIP и двухстадийной диффузионной генерацией. Объединение модальностей повышает качество реконструкции по сравнению с каждой из них в отдельности, а априорная диффузионная модель необходима для семантической согласованности результата.
    \item Разработан метод восстановления отсутствующих каналов ЭЭГ по геометрии монтажа и его вероятностное обобщение на основе маскированного согласования потоков; доказано, что при точном векторном поле обобщение восстанавливает условное распределение пропусков.
\end{enumerate}

Теоретическая значимость результатов заключается в строгом обосновании широко применяемой конструкции касательных признаков: получены утверждения с явными константами, связывающими число каналов, длину ряда и ошибку признаков. Практическая значимость состоит в возможности декодирования по записям одного пациента объемом в десятки объектов, в межсубъектной нейросетевой реконструкции стимула по одновременным записям фМРТ и ЭЭГ и в применимости методов к записям с неполным набором каналов; программные реализации методов опубликованы.

Результаты работы рекомендуется использовать при декодировании по записям одного пациента, когда объем размеченной выборки измеряется десятками объектов и нейросетевые модели неприменимы, для чего предназначены методы оценки лага и классификации, и при построении межсубъектных нейросетевых моделей реконструкции по одновременным записям фМРТ и ЭЭГ, в том числе на записях с неполным набором каналов.

К перспективам дальнейшего развития темы относятся проверка полученных оценок искажения и шума на реальных записях фМРТ и ЭЭГ, обобщение постановки на большее число модальностей с неизвестными взаимными лагами, экспериментальная проверка вероятностного восстановления каналов и его включение в контур декодирования, а также перенос методов на мультимодальные временные ряды иной природы.

Автор выражает благодарность научному руководителю к.ф.-м.н. А.\,В.~Грабовому за постановку задач, обсуждение результатов и научное руководство.


\pdfbookmark{Литература}{bibliography}                                % Закладка pdf

\ifdefmacro{\microtypesetup}{\microtypesetup{protrusion=false}}{} % не рекомендуется применять пакет микротипографики к автоматически генерируемому списку литературы
\urlstyle{rm}                               % ссылки URL обычным шрифтом
\ifnumequal{\value{bibliosel}}{0}{% Встроенная реализация с загрузкой файла через движок bibtex8
    \renewcommand{\bibname}{\large \bibtitleauthor}
    \nocite{*}
    \insertbiblioauthor           % Подключаем Bib-базы
    %\insertbiblioexternal   % !!! bibtex не умеет работать с несколькими библиографиями !!!
}{% Реализация пакетом biblatex через движок biber
    % Цитирования.
    %  * Порядок перечисления определяет порядок в библиографии (только внутри подраздела, если `\insertbiblioauthorgrouped`).
    %  * Если не соблюдать порядок "как для \printbibliography", нумерация в `\insertbiblioauthor` будет кривой.
    %  * Если цитировать каждый источник отдельной командой "--- найти некоторые ошибки будет проще.
    %
    % \nocite{}
    \nocite{dorin2024forecasting}
    \nocite{dorin2025enhancing}
    \nocite{dorin2026decoding}
    \nocite{dorin2025mmro}
    \nocite{dorin2024conf66mipt}
    \nocite{dorin2025conf67mipt}
    \nocite{dorin2026conf68mipt}
    \nocite{dorin2025aij}
    \nocite{dorin2026ioi}

    \ifnumgreater{\value{usefootcite}}{0}{
        \begin{refcontext}[labelprefix={}]
            \ifnum \value{bibgrouped}>0
                \insertbiblioauthorgrouped    % Вывод всех работ автора, сгруппированных по источникам
            \else
                \insertbiblioauthor      % Вывод всех работ автора
            \fi
        \end{refcontext}
    }{
        \ifnum \totvalue{citeexternal}>0
            \begin{refcontext}[labelprefix=A]
                \ifnum \value{bibgrouped}>0
                    \insertbiblioauthorgrouped    % Вывод всех работ автора, сгруппированных по источникам
                \else
                    \insertbiblioauthor      % Вывод всех работ автора
                \fi
            \end{refcontext}
        \else
            \ifnum \value{bibgrouped}>0
                \insertbiblioauthorgrouped    % Вывод всех работ автора, сгруппированных по источникам
            \else
                \insertbiblioauthor      % Вывод всех работ автора
            \fi
        \fi
         % \insertbiblioauthorimportant  % Вывод наиболее значимых работ автора (определяется в файле characteristic во второй section)
        \begin{refcontext}[labelprefix={}]
            \insertbiblioexternal            % Вывод списка литературы, на которую ссылались в тексте автореферата
        \end{refcontext}
        % Невидимый библиографический список для подсчёта количества внешних публикаций
        % Используется, чтобы убрать приставку "А" у работ автора, если в автореферате нет
        % цитирований внешних источников.
        \printbibliography[heading=nobibheading, section=0, env=countexternal, keyword=biblioexternal, resetnumbers=true]%
    }
}
\ifdefmacro{\microtypesetup}{\microtypesetup{protrusion=true}}{}
\urlstyle{tt}                               % возвращаем установки шрифта ссылок URL
